\documentclass[11pt,a4paper]{article}
\usepackage[T1]{fontenc}
\usepackage[utf8]{inputenc}
\usepackage{amsmath,amsthm,mathtools}
\usepackage{newpxtext,newpxmath}
\usepackage[a4paper,margin=25mm,headheight=14pt]{geometry}
\usepackage{microtype}
\usepackage{booktabs,longtable,array,enumitem}
\usepackage{xcolor,listings,fancyhdr,titlesec}
\usepackage{xurl}
\usepackage[hidelinks,pdfencoding=auto]{hyperref}
\definecolor{ink}{HTML}{343E31}
\definecolor{wash}{HTML}{F4F5EF}
\definecolor{rulegray}{HTML}{D5D9CE}
\titleformat{\section}{\Large\bfseries\color{ink}}{\thesection}{0.7em}{}
\titleformat{\subsection}{\large\bfseries\color{ink}}{\thesubsection}{0.7em}{}
\setlength{\parindent}{0pt}
\setlength{\parskip}{5.5pt plus 1pt}
\setlength{\emergencystretch}{3em}
\setcounter{tocdepth}{2}
\setlist{nosep,leftmargin=1.5em}
\pagestyle{fancy}
\fancyhf{}
\fancyhead[L]{\small AsymptoticAnalysis: current-snapshot audit}
\fancyhead[R]{\small \texttt{6687962}}
\fancyfoot[C]{\thepage}
\renewcommand{\headrulewidth}{0.3pt}
\lstset{basicstyle=\ttfamily\small,breaklines=true,breakatwhitespace=false,
 columns=fullflexible,keepspaces=true,showstringspaces=false,
 frame=single,rulecolor=\color{rulegray},backgroundcolor=\color{wash},
 aboveskip=8pt,belowskip=8pt,xleftmargin=3pt,xrightmargin=3pt}
\newcommand{\wl}[1]{\texttt{\expandafter\nolinkurl\expandafter{\detokenize{#1}}}}
\DeclareMathOperator{\Log}{Log}
\newcommand{\OO}{\mathcal{O}}
\newcommand{\RR}{\mathbb{R}}
\newcommand{\CC}{\mathbb{C}}
\newcommand{\id}[1]{\textbf{#1}}
\newtheorem{proposition}{Proposition}[section]
\newtheorem{lemma}[proposition]{Lemma}
\hypersetup{pdftitle={AsymptoticAnalysis: a current-snapshot technical audit},
 pdfauthor={Technical audit prepared for Vladimir Reshetnikov},
 pdfsubject={Source-to-model correctness, Wolfram compatibility, and standalone distribution}}
\begin{document}
\begin{titlepage}
\vspace*{13mm}
{\small\scshape Technical audit \quad / \quad 9 September 2026}\par
\vspace{12mm}
{\Huge\bfseries\color{ink} AsymptoticAnalysis}\par
\vspace{4mm}
{\LARGE A current-snapshot technical audit}\par
\vspace{9mm}
{\large Source-to-model correctness, native Wolfram compatibility,\par
and standalone distribution}\par
\vspace{15mm}
\rule{\textwidth}{0.5pt}
\vspace{4mm}
\textbf{Repository}\par
\url{https://github.com/VladimirReshetnikov/Asymptotic}\par
\vspace{3mm}
\textbf{Pinned revision}\par
\texttt{6687962f3c858a4f93623cfc496f33e35c6763d4}\par
\vspace{3mm}
\textbf{Exclusion baseline}\par
The consolidated issue register and the finding inventories of all eighteen
reports in \wl{external-reports/code-review}.\par
\vspace{8mm}
\textbf{Principal result}\par
A branch-erasing logarithm rewrite can turn a real, eventually invertible source
into the wrong power--log model. Additional findings concern ignored backend
defaults, the standalone builder's dependency gate, and native-result diagnostics.
\vfill
\begin{minipage}{0.97\textwidth}
\small
\textbf{Evidence boundary.} Eleven independent Python tests passed, including
execution of the exact retrieved Python builder on isolated source fixtures.
The mathematical counterexample and its constructive alternative are proved in
this article. Native Wolfram execution was unavailable: public Wolfram
reproducers, nine regression specifications, and the Wolfram portions of the
proposed patches remain unexecuted. No complete upstream-suite pass or native
performance measurement is claimed.
\end{minipage}
\end{titlepage}
\tableofcontents
\newpage

\section{Executive assessment}

The repository is best understood as an asymptotic \emph{modeling and contract}
layer, not merely a collection of extra series coefficients. Its specialized
inverse engines, explicit positive local coordinates, sparse power--log supports,
retained exact cores, source provenance, and selected interval certificates solve
problems different from simply printing a Taylor polynomial. The current native
bridge further separates package analytic objects from results delegated to
\wl{System`Series} or \wl{System`Asymptotic}. That separation is an important
improvement over treating every returned expression as an object with the same
analytic guarantees.\cite{readme,core,native,nativedoc}

The most important new concern is earlier in the pipeline than arithmetic on
remainders. At the symbolic-depth parser's entry, the source expression is
rewritten using
\[
 \log(u^k)\longmapsto k\log u
\]
without proving that $k$ is real. A positive $u$ does not make this identity valid
for a complex exponent. Worse, a subsequent even power can make all resulting
coefficients real, so the recently strengthened real-coefficient checks do not
repair the lost branch information. Section~\ref{sec:winding} gives an explicit
real source, proves that its inverse exists near zero, and shows that the
source-derived erroneous remainder fails by an exponentially growing factor.
The defect is not a request to add general complex asymptotics. It is a request
not to silently erase a branch while claiming a real asymptotic model.\cite{core}

The other principal findings are narrower but independently actionable.

\begin{longtable}{@{}p{0.08\textwidth}p{0.14\textwidth}p{0.43\textwidth}p{0.27\textwidth}@{}}
\toprule
ID & Priority & Finding & Evidence and repair status\\
\midrule
\endhead
N01 & High & Unguarded logarithm-of-power rewriting in symbolic-depth inversion. & Source proof and independent mathematics; two guard edits proposed; native execution pending.\\[4pt]
N02 & Medium & Omitted \wl{"Backend"} is hard-coded as \wl{Automatic}, ignoring \wl{SetOptions}. & Complete dispatcher trace; canonical default-resolution edit proposed; native execution pending.\\[4pt]
N03 & Medium & The standalone dependency gate misses non-bracket application syntax. & Exact retrieved builder executed: eight missed forms; proposed gate rejects all ten dependency forms tested.\\[4pt]
N04 & Low / medium & \wl{"NativeEvaluationStatus"} conflates held user data with unresolved native computation. & Classifier inspected; public characterization pending; diagnostic redesign, not a numerical-error claim.\\
\bottomrule
\end{longtable}

These priorities measure the consequence of the identified mechanism, not its
observed incidence in user programs. N01 is a mathematical identity failure in an
internal normalization path. N02 is a public option-protocol failure. N03 is a
build-time invariant failure under source mutations; it does \emph{not} establish
that the current standalone distribution contains an external dependency. N04 is
a deliberately lower-confidence public diagnostic concern.\cite{core,native,builder}

\subsection{What this audit does not claim}

The Wolfram context connection failed, and a minimal native evaluator request
returned HTTP 502. Neither event is evidence of a package failure. There is no
native kernel version, package-output transcript, or native benchmark from this
audit. Supplied Wolfram expressions are explicitly \emph{reproduction candidates
or desired-contract tests}, not copied observations.

Source access used immutable GitHub connector reads. The Python builder was
materialized exactly from the returned text and matched its returned Git blob
identity. Other package sources were inspected through selected windows or
complete small modules, not obtained as a complete executable checkout. The
source-coverage appendix distinguishes detailed inspection from architectural
review. This article is not a claim that every specialized algorithm or proof in
the repository was independently revalidated.

\subsection{Non-duplication method}

The exclusion baseline is the current consolidated register plus the finding
inventories, README finding summaries, or novelty ledgers of all eighteen
original reports. This is broader than looking at the old top-level
\wl{code-review} location or assuming an earlier snapshot remains current.
The indexed findings are tracked in the supplied machine-readable novelty
ledger.\cite{status,index}

In particular, this report does not present the previously recorded dense native
export, logarithmic remainder loss, generic refinement-demand planning,
coefficient-equality grouping, parameter-diagonal composition, certificate
precision stagnation, or generic native-CI recommendations as fresh discoveries.
The new logarithm issue is compared explicitly with C07 and C16; the option issue
with D01, C09, and the native-compatibility requirement B01. The distinction is
mechanistic, not merely a different example name.

The source now contains fixes discussed by several older reports. For example,
the interval-certificate controller has explicit precision escalation logic,
and the native result contract refuses ordinary analytic operations on a
\wl{"Kind" -> "Native"} object. These are credited as current implementation
facts rather than re-reported as missing features. They are not independently
certified by a native test run here.\cite{cert,native,status}

\section{Architecture and actual capability boundaries}

\subsection{The three layers that must remain distinct}

A useful implementation model has three layers.

\textbf{The request/evaluation layer} receives held Wolfram syntax, distinguishes
package requests from native requests, resolves branches and assumptions, and
selects a backend. \wl{NativeCompatibility.wl} deliberately scans option
containers rather than interpreting arbitrary rules inside the user's source as
wrapper options. It also records held native requests. The source is not just an
already-evaluated algebraic tree: delayed options, held expressions, function
bodies, and evaluation context can matter.\cite{native,inverseexpr}

\textbf{The asymptotic model layer} chooses a positive local coordinate and
constructs sparse coefficient blocks. For an ordinary chart this amounts to a
finite expression such as
\begin{equation}\label{eq:ordinarymodel}
 s(u)=b+A(u)\sum_{\alpha\in E}u^\alpha P_\alpha(\log u)
      +\OO\!\left(|A(u)|u^\rho(1+|\log u|)^d\right),
 \qquad u\downarrow0,
\end{equation}
with an exact offset $b$, a retained carrier $A$, and a finite support $E$.
Different specialized families attach different coefficient algebras or exact
cores to this idea. The shared object head does not, by itself, mean that all
families are closed under every operation.\cite{core,arith,gamma,nativespecial}

\textbf{The evidence layer} records which claims follow from the model and
which additionally concern the original source. A formal native order, a
Poincar\'e remainder, a first-neglected-term bound, a numerical consistency check,
and an interval existence/uniqueness certificate are different evidence objects.
The implementation already reflects parts of this distinction. A complete audit
must not silently replace one with another.\cite{native,dirichlet,special,cert}

N01 attacks the transition from the first layer to the second: once the source
has been changed by an invalid identity, correct coefficient algebra on the
changed model cannot recover the original function. N02 attacks request
resolution. N03 attacks the distribution pipeline that packages these layers.
N04 concerns the description of evidence returned by native delegation.

\subsection{Ordinary inversion and exact supports}

The ordinary inverse constructor identifies a leading term and relative
corrections of the form
\begin{equation}\label{eq:forwardmodel}
 v=a u^p\left(1+\sum_{j=1}^{m}u^{d_j}Q_j(\log u)\right),
 \qquad d_j>0.
\end{equation}
It supports direct Lagrange coefficients, a grouped route, and a Newton route.
For a multi-index $\mathbf{k}$, the correction weight is
$w=\sum_j k_jd_j$. The code implements the coefficient operator, with
$n=\sum_j k_j$, through
\begin{equation}\label{eq:lagrange}
 C_{\mathbf{k}}(L)=
 \frac{(-1)^n r}{p^n\prod_j k_j!}
 \prod_{q=1}^{n-1}\left(\frac{d}{dL}+r+w+pq\right)
 \prod_j Q_j(L)^{k_j}, \qquad n>0.
\end{equation}
For $n=0$ the coefficient is $1$. This is an implementation-relevant advantage:
the package exposes coefficients indexed by exact gaps rather than requiring
all exponents to fit one densely materialized rational grid. This discussion
credits existing machinery; it is not a renewed claim about the already-known
frontier-enumeration bottleneck.\cite{core}

The source distinguishes exponent truncation from marker-depth truncation.
Exponent truncation orders actual asymptotic powers. Depth truncation retains a
finite multi-index simplex and can admit symbolic exponents whose ordering is
proved through assumptions. Those meanings must not be conflated in comparisons
with a native integer order. The N01 witness deliberately selects the latter
mode, because its dedicated finite parser contains the unsafe rewrite.\cite{core}

\subsection{Exact cores and special-function adapters}

The repository exposes logarithmic/Lambert, source-coordinate, retained-core,
Gamma/Barnes, flat-sector, and Fourier-oriented extensions. Their common purpose
is not to expand every nested expression immediately: an appropriate exact
outer transformation or leading inverse can be retained, and corrections are
computed in a better-conditioned coordinate. The inspected Gamma forward path,
for example, combines products in the logarithmic domain and distinguishes
exact recurrence simplification from cancellation in a finite Stirling model.
The special inverse adapters record both the transformed forward model and
its relation to the original function.\cite{readme,modulemap,gamma,special}

The Erfc adapter works with a logarithmic target coordinate and finite tail
models, including a first-neglected-term bound on the positive real tail. The
Gamma adapter similarly records a finite Stirling model rather than claiming
convergence of that asymptotic series. Lambert and quadratic thresholds use
local ramification and an explicit branch choice. These are substantially more
informative contracts than a finite expression alone, provided the source-to-model
identities and domain hypotheses are valid.\cite{special}

The Zeta and Lerch constructions are a different kind of specialization: they
use defining convergent expansions in suitable real regimes and attach explicit
value bounds. Native Wolfram also computes overlapping examples; a list of
function names therefore does not establish a capability exclusive to this
package. The distinction to inspect is the returned representation, bound,
coordinate, and reusable provenance. An earlier review already demonstrated
native Zeta overlap, so that observation is not counted again here.\cite{dirichlet,r17}

\subsection{Native special functions versus an opaque native result}

There are two distinct uses of native Wolfram computation in the current source.
The structured special-function importer calls native \wl{Series} with
\wl{Analytic -> False}, traverses returned \wl{SeriesData} trees, separates
carriers and amplitudes, establishes a real source domain, and transports
absolute errors through supported operations. That is an attempt to import
native coefficient information into a package analytic model.\cite{nativespecial}

By contrast, an explicitly delegated native result is wrapped as
\wl{"Kind" -> "Native"}. Its complete native result is retained; a package
analytic remainder is not fabricated. \wl{requireAnalyticSeries} rejects using
that object as though it possessed the ordinary analytic contract. The user may
still obtain the native result and perform native formal operations on it.
This is a valuable contract barrier.\cite{native,nativedoc}

The first route is an admitted, domain-sensitive importer. The second route is
an interoperability envelope. Calling both ``native support'' without explaining
the distinction makes the API look more uniform than its mathematical contracts
actually are. The article's comparison therefore reports them separately.

\section{Comparison with built-in Wolfram Language}

\subsection{Compare the returned contract, not the spelling of the request}

Current Wolfram documentation describes \wl{Series} as supporting Taylor,
negative and fractional powers, some logarithmic expansions, expansion at
infinity, and successive expansion in several variables. It usually returns a
\wl{SeriesData} object. It is not accurate to characterize it as an
integer-power-only Taylor tool.\cite{wlseries}

\wl{SeriesData} encodes a rational power lattice and permits logarithmic or
slower-growing expressions in coefficients; symbolic powers and rapidly growing
or oscillating factors may remain outside the object. Native arithmetic,
integration, differentiation, composition, and inversion act on that formal
representation. This is not the same encoding as a package pair $(\rho,d)$
that explicitly records a logarithmic degree in an analytic remainder.\cite{wlseriesdata}

\wl{Asymptotic} is broader than a request for a rational-lattice series. Its
documentation permits automatically inferred scales, complex or infinite
centers, and expressions involving integrals, transforms, or differential-equation
solutions. Its approximation order is not generally a polynomial degree.
A finite native asymptotic answer must be interpreted using that function's
contract, not by attaching an invented package \wl{PowerLogRemainder}.
\cite{wlasymptotic}

\begin{longtable}{@{}p{0.20\textwidth}p{0.35\textwidth}p{0.37\textwidth}@{}}
\toprule
Task & Built-in Wolfram baseline & Current package role and limitation\\
\midrule
\endhead
Ordinary local series & \wl{Series}; native coefficient and formal-order arithmetic. & Adds its own explicit positive chart and remainder metadata; ordinary overlap is substantial.\\[5pt]
Power--log, irrational-gap inversion & \wl{InverseSeries} for native series; \wl{AsymptoticSolve} for implicit asymptotics. & Dedicated exact-gap and logarithmic coefficient machinery; no blanket claim that native tools fail on every overlapping expression.\\[5pt]
General implicit systems & \wl{AsymptoticSolve} supports systems and multivariate independent variables. & Specialized scalar inverse machinery is not an independently implemented general system solver.\\[5pt]
Real branch selection & Native solvers accept real domains, assumptions, and directional specifications. & Positive source charts, source provenance, and explicit inverse-family selection are central package contracts.\\[5pt]
Gamma/Erfc and related tails & Built-in functions, \wl{Series}, \wl{Asymptotic}, and implicit solution tools overlap many forward cases. & Retained inverse cores, finite-model provenance, and selected transported error bounds are the meaningful additions.\\[5pt]
Many other special functions & Native coefficient/asymptotic algorithms are a principal computational source. & Structured importer adds admissibility and envelope handling; opaque delegation adds coverage but not an analytic proof.\\[5pt]
Complex and multivariate expansions & Documented native scope is broad. & May be delegated under a native result contract; this does not imply a general complex/multivariate package remainder algebra.\\[5pt]
Composition & \wl{ComposeSeries} requires matching outer center and inner limiting value. & Adds specialized chart and evidence handling, but shared object syntax does not imply universal closure.\\[5pt]
Integrals and differential equations & \wl{AsymptoticIntegrate} and \wl{AsymptoticDSolveValue} are dedicated engines. & Not independently reimplemented by the inspected scalar inverse layer. Native delegation is not equivalent to an imported analytic certificate.\\[5pt]
Numerical inverse checking & \wl{FindRoot} and arbitrary-precision evaluation offer numerical evidence. & Dedicated checks preserve source intent; supported interval certificates additionally establish bounded claims.\\
\bottomrule
\end{longtable}
The comparison is grounded in the current package sources and the documented
native interfaces; it is not a fresh empirical success/failure matrix.\cite{core,native,nativedoc,special,cert,wlinverseseries,wlasymsolve,wlcompose,wlasymintegrate,wlasymdsolve}

\subsection{Three concrete comparison patterns}

\textbf{An aligned ordinary test.} For $e^x$, a package exclusive cutoff $4$
and a native inclusive polynomial order $3$ both retain
\[
 1+x+\frac{x^2}{2}+\frac{x^3}{6}.
\]
The comparison script uses that alignment explicitly:
\begin{lstlisting}
AsymptoticExpansion[Exp[x], {x, 0, 4}, "Backend" -> "Package"]
System`Series[Exp[x], {x, 0, 3}]
\end{lstlisting}
The coefficient identity follows directly from the exponential series. It is
not reported as native output observed during this audit. The request semantics
come from the two interfaces.\cite{core,wlseries}

\textbf{A logarithmic test.} For $\exp(x\log x)$, the same alignment compares
retained algebraic powers, but the native displayed order and the package's
log-degree metadata need separate examination. Comparing only
\wl{Normal[result]} would discard precisely the information being audited.
The comparison script records both the finite expression and remainder-related
fields, rather than treating equality after \wl{Normal} as sufficient.
This follows from the representations' definitions, not from a newly asserted
failure in the old export issue.\cite{core,wlseriesdata}

\textbf{An irrational-gap inverse.} For
\[
 y=x+x^{1+\sqrt2},\qquad x\downarrow0,
\]
the dedicated inverse route and a real \wl{AsymptoticSolve} request address the
same source equation. Their order arguments need not mean the same thing.
One should compare actual exponents, the selected root, the residual against the
original equation, and the remainder meaning. The supplied requests are
coverage probes; no result or speed ranking is claimed. Native
\wl{AsymptoticSolve} explicitly supports real solutions, implicit functions,
systems, and automatically inferred scales.\cite{core,wlasymsolve}

\subsection{Where the package can be preferable}

The package is particularly attractive when the consumer needs more than a
finite formula: an exact sparse support, a named inverse branch, an exact
logarithmic or Lambert core, explicit input-model truncation, subsequent
precision-aware operations, or a source-linked certificate. Those features
should be assessed by checking that their metadata really describes the
expression and source at hand. They do not become valid merely because the
object carries an \wl{Assumptions} key.\cite{core,special,cert}

A strong implementation strategy is consequently to use native Wolfram as a
coefficient or solution oracle where suitable, while keeping branch admission,
remainder import, source provenance, and supported operations explicit. The
current separation of structured import and opaque native delegation follows
that strategy in part. N01 demonstrates why branch checks must precede
coefficient-domain checks; N02 demonstrates why request defaults must be resolved
before choosing which contract will be returned.\cite{native,nativespecial,core}

\subsection{Where the built-ins remain the natural first tool}

For ordinary formal power-series manipulation, native \wl{InverseSeries} and
\wl{ComposeSeries} are direct tools. The former inverts an appropriate native
series; the latter composes series with compatible centers. Native
\wl{SeriesData} arithmetic already propagates its formal truncation order.
A wrapper is not automatically a better replacement for those operations.
\cite{wlinverseseries,wlcompose,wlseriesdata}

For asymptotic integrals or differential equations, Wolfram exposes dedicated
engines rather than requiring a prior scalar inverse reduction.
\wl{AsymptoticIntegrate} accepts integral problems; \wl{AsymptoticDSolveValue}
accepts equations, systems, and parameter-asymptotic forms. These capabilities
are relevant comparison baselines, not a demand that the package independently
reimplement them. The broad wishlist is already in the older reports and is not
repeated as a new development finding.\cite{wlasymintegrate,wlasymdsolve,status}

\section{N01: branch erasure in the symbolic-depth parser}\label{sec:winding}

\subsection{The exact source mechanism}

The relevant source is \wl{parseFinite} in
\wl{src/Kernel/AsymptoticAnalysis.wl}, lines 916--928. Line 917 contains two
rewrite rules; formatted separately, they are:
\begin{lstlisting}
Log[u^k_] /; FreeQ[k, u] :> k Log[u]

Log[c_ u^k_.] /;
  FreeQ[c, u] && FreeQ[k, u] &&
  TrueQ[Simplify[c > 0, ass]] :> Log[c] + k Log[u]
\end{lstlisting}
The exponent is required to be independent of $u$, but not real. The second
rule proves positivity of the multiplicative coefficient $c$ but again does
not establish realness of $k$. These rules execute before the logarithm is
replaced by the formal polynomial variable and before
\wl{rowsToModel} checks real coefficients.\cite{core}

For a positive real base, the standard identity is safe for a real exponent.
For $k=a+ib$, however,
\begin{equation}\label{eq:principalpower}
 \Log(u^k)=a\log u+i\operatorname{Arg}\!\left(e^{ib\log u}\right),
 \qquad u>0,
\end{equation}
whereas the rewrite produces $a\log u+ib\log u$.
The difference is an integer multiple of $2\pi i$, generally depending on $u$.
Wolfram's documentation explicitly warns that unrestricted power expansion
ignores branch information and need not preserve numerical values. The package
rule performs one of those branch-sensitive transformations without its missing
hypothesis.\cite{wllog,wlpowerexpand}

\begin{proposition}\label{prop:noeventual}
For a fixed $k\in\CC$ with $\operatorname{Im}k\ne0$, the identity
$\Log(u^k)=k\log u$ does not hold throughout any punctured positive interval
$(0,\varepsilon)$.
\end{proposition}
\begin{proof}
Write $b=\operatorname{Im}k$. The imaginary part of the principal logarithm in
\eqref{eq:principalpower} remains in a bounded interval of length $2\pi$.
The imaginary part of $k\log u$ is $b\log u$, which is unbounded as
$u\downarrow0$. Equality on an entire sufficiently small positive interval is
therefore impossible.
\end{proof}

This is stronger than exhibiting a branch-cut exception far from the expansion
point. No amount of restricting the positive source variable to a smaller
neighborhood makes this particular rewrite eventually valid when
$\operatorname{Im}k\ne0$.

\subsection{Why real-coefficient validation does not save the model}

Consider the parameterized input
\begin{equation}\label{eq:badsource}
 f(x)=x+x^2\Log(x^a)^2,
 \qquad a^2=-1,
 \qquad x>0.
\end{equation}
The expression contains a symbolic parameter $a$, not a literal \wl{Complex}
atom. Its source syntax passes the core's syntactic exactness condition. After
the unsafe rewrite it becomes
\[
 \widetilde f(x)=x+a^2x^2(\log x)^2
               =x-x^2(\log x)^2.
\]
All coefficients of this changed power--log model are real under the stated
assumption. Thus a gate that validates only the resulting coefficients can
approve the wrong model.\cite{core}

The two possible parameter values $a=i$ and $a=-i$ produce the \emph{same real}
source in \eqref{eq:badsource}. Define
\begin{equation}\label{eq:periodicP}
 \theta(t)=\operatorname{Arg}(e^{it}),\qquad
 P(t)=\theta(t)^2.
\end{equation}
At the two choices of principal endpoint, the square has the same value.
Consequently $P$ is continuous and $2\pi$-periodic, and
\begin{equation}\label{eq:realF}
 f(x)=x-x^2P(\log x),\qquad 0\le P(t)\le\pi^2.
\end{equation}
The actual perturbation is bounded by $\pi^2x^2$. The rewritten perturbation is
$x^2(\log x)^2$, which has a different leading amplitude. This is not a merely
complex-valued input outside a real-only output contract: the source itself is
real on the entire positive approach.

\subsection{Existence of the actual real inverse}

On every interval away from the phase corners,
$P'(t)=2\theta(t)$, so $|P'|\le2\pi$. Continuity across the corners makes $P$
globally Lipschitz with constant $2\pi$. Set
\begin{equation}\label{eq:delta}
 \delta=\frac{1}{4(\pi^2+\pi)}.
\end{equation}
Where $f$ is differentiable,
\[
 f'(x)=1-2x\bigl(\theta(\log x)^2+\theta(\log x)\bigr)
 \ge 1-2x(\pi^2+\pi)\ge\frac12,
 \qquad 0<x\le\delta.
\]
There are finitely many corners on any compact subinterval of $(0,\delta]$.
Continuity and the derivative lower bound on the pieces therefore imply strict
increase on that interval. Also $f(x)\to0$ as $x\downarrow0$.
Thus a unique real inverse $g$ exists on $(0,f(\delta))$.

The source is not asserted to be an analytic power--log germ: its derivative
has phase corners accumulating at zero. That is precisely why a sound parser
must either retain the periodic coefficient in an admitted representation or
reject the unsupported normalization. It must not replace it by an unrelated
analytic-looking model.

\subsection{The inverse returned by the changed model}

For the changed model $\widetilde f(x)=x-x^2(\log x)^2$, the depth-one
Lagrange construction has leading power $p=1$, one gap $d_1=1$, and
$Q_1(L)=-L^2$. Formula~\eqref{eq:lagrange} gives the first correction $L^2$.
The depth remainder formula in \wl{construct} assigns power $3$ and logarithmic
degree $4$. Thus the source-derived core result is
\begin{equation}\label{eq:falseinverse}
 \widetilde g_1(y)=y+y^2(\log y)^2,
 \qquad \OO\!\left(y^3(1+|\log y|)^4\right).
\end{equation}
That is a valid type of approximation for the \emph{changed} source, but not
for \eqref{eq:realF}. This algebraic trace is not a captured public Wolfram
output; the supplied probes separate private normalization, core construction,
and the complete public facade.\cite{core}

Choose the exact sequence
\[
 x_n=y_n=e^{-2\pi n},\qquad n=1,2,\ldots.
\]
Then $x_n^{\pm i}=1$, so $P(\log x_n)=0$, $f(x_n)=x_n$, and
$g(y_n)=y_n$. The error in \eqref{eq:falseinverse} is exactly
$y_n^2(2\pi n)^2$. Its ratio to the claimed remainder scale is
\begin{equation}\label{eq:divergenceratio}
 \frac{|\widetilde g_1(y_n)-g(y_n)|}
      {y_n^3(1+|\log y_n|)^4}
 =\frac{e^{2\pi n}(2\pi n)^2}{(1+2\pi n)^4}
 \longrightarrow\infty.
\end{equation}
This proves failure of the advertised asymptotic scale for the original source;
no fitted constants or numerical root solver are needed for the conclusion.

\begin{center}
\begin{tabular}{@{}rrr@{}}
\toprule
$n$ & $y_n$ & Ratio in \eqref{eq:divergenceratio}\\
\midrule
1 & $1.86744\,10^{-3}$ & $7.51324$\\
2 & $3.48734\,10^{-6}$ & $1.33681\,10^{3}$\\
3 & $6.51241\,10^{-9}$ & $3.51445\,10^{5}$\\
5 & $2.27110\,10^{-14}$ & $3.93577\,10^{10}$\\
10 & $5.15790\,10^{-28}$ & $4.61038\,10^{23}$\\
20 & $2.66039\,10^{-55}$ & $2.30603\,10^{50}$\\
\bottomrule
\end{tabular}
\end{center}
The displayed values were independently evaluated with mpmath. They illustrate
the exact formula, rather than serving as a substitute for its proof.

\subsection{Reproduction candidates and the smallest safe repair}

The public candidate is:
\begin{lstlisting}
Clear[a, x, y];
AsymptoticInverse[
  x + x^2 Log[x^a]^2,
  {x, 0}, {y, 1},
  Assumptions -> a^2 == -1,
  "Truncation" -> "Depth"]
\end{lstlisting}
For localization, the characterization script additionally calls
\wl{parseFinite} directly and invokes the core constructor under the package's
failure catcher. These distinguish an identity defect in the private parser
from any earlier public-family routing or rejection. A future native run should
record all three, rather than declaring the public result solely from this
source trace.

Both rewrite rules should require
\begin{lstlisting}
TrueQ[Simplify[Element[k, Reals], ass]]
\end{lstlisting}
in addition to their existing conditions. For a fixed real $k$ and positive
$u$, $u^k$ is positive, and the identity is valid. In the scaled rule, the existing
proof $c>0$ completes the sufficient conditions. An unproved condition should
leave the logarithm unnormalized; this finite power--log parser will then reject
it rather than fabricate coefficients.

This is an intentionally conservative patch. It does not prove that every
normalization in every other engine is sound, and it does not introduce a
complex-asymptotic mode. It also does not convert an inconclusive proof into
``false'': it simply refuses to use the identity without a sufficient proof.
The proposed edit tool patches exactly these two anchors, and the Wolfram tests
check both the rejection case and a real-exponent positive control.

\subsection{Novelty relative to C07 and C16}

C07 concerns admission of nonreal coefficients into a purportedly real result.
Here the coefficients \emph{after rewriting} really are real. The failed
obligation is equality between the source and the normalized expression.
C16 concerns treating formal Taylor information for an opaque source as analytic
information. Here the source is explicit and the failure is a specific branch
identity, not an unknown derivative or unjustified smoothness assumption.
Neither older repair, by itself, addresses this mechanism.\cite{status,r11,r18}

The appropriate invariant is therefore a chain:
\[
 \text{valid source identity}
 \quad\Longrightarrow\quad
 \text{admitted coefficient model}
 \quad\Longrightarrow\quad
 \text{valid remainder algebra}.
\]
A check at the second or third stage cannot compensate for a failed implication
at the first.

\section{A constructive alternative: retain the winding coefficient}\label{sec:constructive}

Rejecting the unsafe rewrite is the immediate correctness repair. The same
example also identifies a concrete extension beyond rejection: retain the bounded
periodic function $P(\log y)$ instead of expanding it as an unbounded polynomial
in $\log y$.

\begin{proposition}\label{prop:goodinverse}
For the real source \eqref{eq:realF}, its inverse satisfies
\begin{equation}\label{eq:goodinverse}
 g(y)=y+y^2P(\log y)+\OO(y^3),\qquad y\downarrow0.
\end{equation}
More explicitly, for $0<y<3\delta/4$,
\begin{equation}\label{eq:explicitbound}
 \left|g(y)-\left(y+y^2P(\log y)\right)\right|
 \le K y^3,
 \qquad
 K=\frac{112}{27}\pi^4+\frac{32}{9}\pi^3
   \approx514.311879893224285,
\end{equation}
where $\delta$ is given by \eqref{eq:delta}.
\end{proposition}
\begin{proof}
Let $x=g(y)$. Since $x\le\delta$ and $P\le\pi^2$,
\[
 y=x-x^2P(\log x)\ge\frac34x.
\]
The assumed bound on $y$ ensures $x<\delta$, because
$f(\delta)\ge3\delta/4$. Therefore
\[
 y\le x\le\frac43y,
 \qquad 0\le x-y\le\pi^2x^2\le\frac{16}{9}\pi^2y^2.
\]
Using the $2\pi$ Lipschitz constant of $P$,
\begin{align*}
 |x-y-y^2P(\log y)|
 &=|x^2P(\log x)-y^2P(\log y)|\\
 &\le\pi^2(x-y)(x+y)+2\pi y^2\log(x/y)\\
 &\le\pi^2\left(\frac{16}{9}\pi^2y^2\right)
              \left(\frac73y\right)
      +2\pi y^2\left(\frac{16}{9}\pi^2y\right),
\end{align*}
where $\log(x/y)\le(x-y)/y$ was used. The resulting constant is exactly $K$.
\end{proof}

The independent checks evaluated this bound at sixty targets with high-precision
bisection. Those tests are useful implementation sanity checks, but the
uniform analytic conclusion is the proposition, not the floating-point sample.
A separate Wolfram reference file returns the approximation and its conditional
bound as ordinary expressions. It does not claim to install a new package
coefficient algebra or to generate an \wl{InverseCertificate}.

\subsection{Why a fixed finite Fourier approximation is insufficient}

On $[-\pi,\pi]$, the coefficient is $P(t)=t^2$, periodically continued. Direct
integration of its Fourier coefficients gives
\begin{equation}\label{eq:fourierP}
 P(t)=\frac{\pi^2}{3}
      +4\sum_{m=1}^{\infty}\frac{(-1)^m\cos(mt)}{m^2}.
\end{equation}
The series is absolutely and uniformly convergent. If $P_M$ is its ordinary
partial sum, then
\[
 \|P-P_M\|_\infty
 =4\sum_{m>M}\frac1{m^2}\asymp\frac4M.
\]
The upper bound follows by the triangle inequality, and equality is attained
at $t=\pi$, where the terms all have the same sign. Replacing $P$ by $P_M$ in
\eqref{eq:goodinverse} introduces an error of order $y^2/M$. Consequently, to
preserve a \emph{uniform} cubic bound using this partial-sum strategy, one needs
$M$ of order at least $1/y$. A fixed number of harmonics does not preserve the
next asymptotic order.

This is a specific representation-cost argument, not a benchmark of the
repository's Fourier implementation and not a lower bound for every possible
representation of $P$. In fact, retaining $P$ exactly as a periodic piecewise
quadratic coefficient has constant structural size. It provides a concrete
reason to support a bounded coefficient with a Lipschitz contract instead of
forcing every such coefficient into a fixed finite harmonic polynomial.

\subsection{A narrowly scoped implementation opportunity}

A prototype coefficient object could record a held function, its period, a
supremum bound, a Lipschitz constant, and the domain of its parameters. For this
example the data are
\[
 (P,\;2\pi,\;\pi^2,\;2\pi).
\]
The first supported operation need only be the one-step inverse estimate used
in Proposition~\ref{prop:goodinverse}. It requires a value bound and a Lipschitz
bound, not analyticity or an unrestricted derivative series. This is narrower
and more testable than announcing a general transseries extension.

For nonnegative bounded Lipschitz $P$ with $\|P\|_\infty\le M$ and Lipschitz
constant $L$, where $M+L>0$, the same proof works after taking
$\delta\le1/[2(2M+L)]$. Its cubic error constant can be taken as
\[
 \frac{112}{27}M^2+\frac{16}{9}ML.
\]
A useful implementation milestone is therefore an explicit, small theorem-backed
constructor for $x-x^2P(\log x)$, with exact coefficient retention and no claimed
higher derivative regularity. General closure, higher-order corner behavior,
and integration remain separate tasks. This proposal is derived from N01, not
an unimplemented feature asserted to exist today.

\section{N02: declared backend defaults are not consumed}\label{sec:defaults}

\subsection{The mismatch}

\wl{NativeCompatibility.wl} merges a \wl{"Backend" -> Automatic} entry and
native option lists into \wl{Options[AsymptoticExpansion]}. It also copies those
options onto the short alias \wl{AsymptoticExpand}. The dispatcher then collects
explicit selector rules from held trailing option containers. In the absence of
such a rule, it does not consult the declared defaults:
\begin{lstlisting}
backend = If[values === {}, Automatic, ReleaseHold[First[values]]];
\end{lstlisting}
Thus the option table and the dispatch implementation are disconnected for this
key.\cite{native}

Wolfram's \wl{SetOptions} interface changes the default option settings associated
with a symbol. A function that advertises configurable options normally obtains
its effective values through that option protocol. The current dispatcher
instead interprets ``not explicitly supplied'' as ``always Automatic.''
\cite{wlsetoptions}

A source-derived public witness is:
\begin{lstlisting}
saved = Options[AsymptoticExpansion];
Internal`WithLocalSettings[
  SetOptions[AsymptoticExpansion, "Backend" -> "Series"],
  s = AsymptoticExpansion[Exp[x], {x, 0, 3}];
  {s["Kind"], s["NativeBackend"], Normal[s]},
  Options[AsymptoticExpansion] = saved
]
\end{lstlisting}
The expected option-protocol behavior is a native Series result. The inspected
source instead selects \wl{Automatic}; for this ordinary exact request, the
package path is the natural successful route. The resulting order convention
can differ as well: package cutoff $3$ omits the cubic term, whereas native
Series order $3$ retains it. The exact public result remains a native
reproduction candidate, not an observation made here.\cite{core,native,wlseries}

\subsection{The minimal repair and its boundary}

The narrow replacement is:
\begin{lstlisting}
backend = If[values === {},
  OptionValue[AsymptoticExpansion, {}, "Backend"],
  ReleaseHold[First[values]]];
\end{lstlisting}
This preserves the explicit selector handling and resolves the canonical default
only when no explicit selector exists. The test specifications cover a constant
default, an explicit override, a delayed default with a counter, the alias under
a canonical default, and the unchanged explicit-native control.

The patch does not silently invent a policy for every merged option. Native
options on the wrapper, native backend defaults, package defaults, and explicit
package-only protections are separate concerns. For example, blindly passing
\emph{all} defaults into the existing automatic-protection logic could mark
nearly every request as package-only merely because \wl{"MaxTerms"} has a
default value. The dispatcher must know both a value and its origin.
\cite{native,nativedoc}

\subsection{The alias needs an explicit default policy}

The source defines
\begin{lstlisting}
Options[AsymptoticExpand] = Options[AsymptoticExpansion];
AsymptoticExpand[args___] := AsymptoticExpansion[args];
\end{lstlisting}
The alias's own option table is not read by the delegation. The minimal patch
makes a canonical default work through the alias, but changing the alias's
\emph{own} defaults still does not automatically change canonical dispatch.
This should not be left to users to infer from copied \wl{Options} output.
\cite{native}

There are two coherent designs. A true synonym can document and expose one
canonical configuration owner. Alternatively, the alias can be an independent
entry point whose defaults are resolved first and then passed to a shared
implementation. The latter must preserve the distinction between an explicit
caller request and a wrapper-resolved default. Copying an option list once while
ignoring it at runtime is neither design.

\subsection{A precise request-resolution refinement}

The earlier reports already recommend typed requests and clearer cutoff
semantics. This report does not repeat that general recommendation as new.
The concrete delta is to retain the \emph{origin} of each resolved selector:
\begin{lstlisting}
<|"Value" -> "Series",
  "Origin" -> "CanonicalDefault",
  "Owner" -> AsymptoticExpansion,
  "Evaluated" -> True|>
\end{lstlisting}
An explicit selector has origin \wl{"Caller"}; an inherited native default is
another case. Backend selection consumes a resolved value once. Protection
logic can then deliberately inspect caller intent rather than accidentally
using the presence of every effective default as a veto.

A useful invariant is that the returned object records the backend actually
chosen and the effective order meaning. A useful test is that delayed defaults
are not re-evaluated solely for metadata. Neither requirement demands that all
native and package semantics become identical. It demands that dispatch be
explainable and consistent with the advertised option protocol.

\section{N03: the standalone dependency gate is syntax-incomplete}\label{sec:builder}

\subsection{What the builder actually verifies}

The Python builder recursively inlines recognized companion loads into one
standalone Wolfram file. It masks strings and nested comments, checks remaining
executable text for forbidden dependency operations or runtime file-location
symbols, and emits deterministic source boundaries. The dedicated upstream
Python tests inspect load order, missing companions, cycles, inert text,
freshness, and several bracket-form dependencies. These are useful protections.
\cite{builder,buildtests}

The incomplete part is the dependency expression:
\begin{lstlisting}
(?:Get|Needs|Import|OpenRead|ReadList|Read|BinaryRead|
   URLRead|URLExecute|URLDownload)\s*\[
\end{lstlisting}
It recognizes selected heads only when followed by a bracket. Wolfram source
also admits prefix, postfix, and higher-order application forms. Thus
\wl{Get @ "missing.wl"}, \wl{"missing.wl" // Get}, and
\wl{Apply[Get, {"missing.wl"}]} are not recognized by that gate. The preceding
companion-load inliner also does not rewrite these forms.\cite{builder,wlexpressions,wlapply}

\subsection{Executed mutation evidence}

The exact retrieved builder was loaded as a Python module. Each test created a
temporary miniature \wl{src/Kernel} tree with a valid entry-point directory setup,
inserted one candidate statement, and called the builder's actual
\wl{assemble()} function. No Wolfram fixture content was executed and no actual
network request was made.

\begin{longtable}{@{}p{0.63\textwidth}p{0.13\textwidth}p{0.16\textwidth}@{}}
\toprule
Inserted executable form & Baseline & Proposed gate\\
\midrule
\endhead
\wl{Get["missing.wl"]} & Reject & Reject\\
\wl{System`Get["missing.wl"]} & Reject & Reject\\
\wl{Get @ "missing.wl"} & Accept & Reject\\
\wl{"missing.wl" // Get} & Accept & Reject\\
\wl{Import @ "missing.json"} & Accept & Reject\\
\wl{"Missing`" // Needs} & Accept & Reject\\
\wl{Apply[Get, {"missing.wl"}]} & Accept & Reject\\
\wl{System`Get @ "missing.wl"} & Accept & Reject\\
\wl{ReadList @ "missing.wl"} & Accept & Reject\\
\wl{URLDownload @ "https://example.invalid/x"} & Accept & Reject\\
\bottomrule
\end{longtable}

The four controls cover quoted text, nested comments, unrelated larger symbols
including \wl{Get$Helper}, and an ordinary function definition. Both versions
accept them. The proposed version also preserves the intended direct companion
inlining. The complete independent suite contains eleven test methods and
additional subcases; its actual log and structured mutation results are supplied.

This proves a packaging-invariant gap under concrete source edits. It does not
prove that an existing module uses any missed form, and it is not an allegation
of an external load hidden in the current distribution. The observed failure is
that the advertised gate would allow such an accidental dependency to be
introduced unnoticed.

\subsection{The tested conservative repair}

After comment and string masking, the proposed gate recognizes the forbidden
\emph{symbol tokens} regardless of how they are applied. Identifier boundaries
must include \wl{$} as well as ordinary word characters, so names such as
\wl{Get$Helper} are not mistaken for \wl{Get}. A context-qualified occurrence
such as \wl{System`Get} remains detectable.

This deliberately conservative policy can reject a forbidden head stored as
held data, or a same-leaf-name symbol from another context. Such a rejection is
a request for source review, not a claim that the occurrence necessarily executes
I/O. That is an acceptable tradeoff for a small self-contained distribution
builder, provided the policy is stated. The existing gate was already
conservative about bracket-form calls inside held code.\cite{builder}

The repair is not a security sandbox and does not decide arbitrary program
behavior. A dynamically constructed function name or other reflective mechanism
can evade a literal-token check. Proving absence of all possible runtime I/O in
an unrestricted symbolic language is not the builder's present contract. The
specific implementable requirement is to catch ordinary literal dependencies in
all common application syntaxes, and to make unrecognized dynamic evaluation an
explicit review category rather than implying a total proof.

\subsection{Performance and integration}

The existing builder already performs a linear masking pass and a scan of the
resulting source. The proposed token-boundary check preserves that structural
cost. No extra Wolfram kernel is needed to run its mutation regressions. This
is a focused build-pipeline improvement, distinct from the older general
recommendation to add native mathematical CI.\cite{builder,workflow,status}

The patch tool edits only \wl{validation/build_standalone.py}, when N03 is
selected. It does not regenerate the standalone distribution automatically.
That keeps an executable source change separate from artifact production and
allows the existing builder tests plus the new mutation tests to be reviewed
before publishing a new generated file.

\section{N04: native-status classification should describe what it knows}

The native wrapper computes its status using a syntactic search:
\begin{lstlisting}
If[FreeQ[result, _System`Series | _System`Asymptotic],
  "Computed", "Unresolved"]
\end{lstlisting}
This detects occurrences of those heads anywhere in the result tree. It does
not distinguish an unresolved active computation from a native call stored as
held user data.\cite{native}

A minimal public characterization candidate compares native and wrapped expansion
of an expression constant in $x$:
\begin{lstlisting}
Series[HoldComplete[Series[t, {t, 0, 1}]], {x, 0, 1}]

AsymptoticExpansion[
  HoldComplete[Series[t, {t, 0, 1}]], {x, 0, 1},
  "Backend" -> "Series"]
\end{lstlisting}
If the native operation returns the held constant unchanged, the wrapper's
classifier labels it \wl{"Unresolved"} because the \wl{Series} stored inside
\wl{HoldComplete} is present. That stored expression is data, not a failed
attempt by the outer backend to expand with respect to $x$. The exact public
native behavior was not observed here; the script records both raw results to
settle it without assuming the expected native simplification.

The current development document is careful not to equate a computed native
result with a package analytic proof. Nevertheless, the property name invites a
stronger reading than this syntactic predicate supports. The immediate issue is
misleading status or future fallback decisions, not an established numerical
wrong answer.\cite{nativedoc}

A conservative design separates three facts: the backend returned a value;
residual native-call syntax occurs at specified positions; and the system has or
has not established that those occurrences represent active unresolved work.
For opaque held data, \wl{"ReturnedOpaque"} or an explicitly syntactic property
is more accurate than a confident failure classification. A traversal must not
release held bodies merely to classify them. If the implementation cannot safely
interpret an arbitrary head's evaluation behavior, an unknown status is preferable
to executing user data.

No code patch for this diagnostic policy is included. The raw \wl{NativeResult}
should remain unchanged, and the existing analytic/native contract barrier
should remain in force. The supplied characterization is intended to establish
the exact returned shape before choosing an API-compatible status migration.

\section{Focused performance, API, and development conclusions}

\subsection{Measure complete semantic jobs}

A benchmark comparing a package analytic object with \wl{Normal[Series[...]]}
can compare different jobs: branch proof and metadata construction on one side,
coefficient generation with erased order information on the other. Conversely,
measuring only an internal coefficient recurrence can conceal facade and source
normalization costs. The new regression work should therefore separate request
resolution, source admission, coefficient construction, and optional result
views. This is a measurement protocol, not a native speed claim.\cite{core,native}

The supplied comparison runner records a single elapsed time per bounded call,
its messages, head, finite expression representation, and remainder fields.
Those timings are observations to inspect, not published performance rankings.
Only the explicitly aligned ordinary/logarithmic forward pairs have matching
algebraic cutoffs. The inverse and special-function pairs are scope probes;
after a native run, their actual branches and orders must be compared before
any timing summary is meaningful.

The two new concrete cost conclusions are limited and testable. First, the
builder repair retains a linear lexical pass. Second, ordinary finite Fourier
partial sums of the winding coefficient require an order-dependent number of
harmonics to preserve the cubic uniform inverse bound. Neither conclusion
recycles the older general enumeration or provenance-size findings.

\subsection{Source identities need local proof obligations}

The N01 patch is small, but the development lesson is precise. A normalization
rule should declare its sufficient conditions at the point where it changes the
source, rather than relying on the realness of the final polynomial. For the
identified rule those conditions are $u>0$, $k\in\RR$, and, when present, $c>0$.

A practical follow-up is a small registry of branch-sensitive identities used
by source normalizers, each with a positive control, a negative control, and a
boundary/winding family. The first family can vary $a^2=-c^2$ for positive
rational $c$, use even powers of the logarithm so the source remains real, and
sample exact winding points $x_n=e^{-2\pi n/c}$. This is much more targeted than
adding random expression tests and hoping branch effects appear.

Proof caching, if later added for these guards, must key on the expression,
parameter assumptions, and chart. A proof obtained for a real exponent cannot
be reused merely because a later normalized coefficient looks identical. This
is a design consequence of the counterexample, not a claim that the current
repository contains such a cache defect.

\subsection{Backend defaults need one owner and one evaluation}

The first milestone is the N02 repair with its explicit and delayed-default
controls. The second is a documented alias ownership policy. The third is an
origin-aware effective request record that distinguishes caller intent from
inherited defaults. These milestones are small enough to validate independently
of the broad, explicitly unfinished native-superset requirement.\cite{native,nativedoc}

In particular, completion of native fallback coverage should not be used to
postpone a deterministic option-protocol repair. Nor should this repair be sold
as completion of B01: accepting all syntax that either native backend accepts,
selecting a second backend after an unresolved first attempt, and preserving
arbitrary evaluation-context effects are larger obligations already discussed
in the repository's compatibility document.

\subsection{Keep new functionality tied to an explicit theorem}

The bounded periodic inverse of Section~\ref{sec:constructive} is a useful
experimental extension because its assumptions, approximation, and error
constant are explicit. A first implementation can remain outside the main
package and return an ordinary expression plus a conditional bound. The supplied
\wl{WindingAwareReference.wl} does exactly that at the formula level; it has not
been executed in a native kernel.

Only after that reference behavior is validated should it be integrated with a
new coefficient representation. The first package contract should not claim
arbitrary-order differentiability, because the coefficient has corners. It
should claim only the operations justified by its value and Lipschitz bounds.
This avoids recreating the original error under a more general object head.

\section{Validation record and how to use the artifacts}

\subsection{Executed work}

The independent test suite passed eleven tests. Its subcases include the exact
retrieved builder's ten dependency forms, eight baseline misses, the proposed
gate's rejections, inert and unrelated-symbol controls, companion inlining,
patch-anchor refusal behavior, principal-log evaluation, the divergence formula,
Lipschitz samples, and sixty numerical checks of the constructive inverse bound.
The source identity of the local builder fixture was checked against the Git
blob identifier returned by the connector. No checksum file is needed for this
in-memory verification.

The exact mathematical conclusion for N01 is Proposition~\ref{prop:noeventual}
and equation~\eqref{eq:divergenceratio}. The constructive error bound is
Proposition~\ref{prop:goodinverse}. Floating-point checks exercise formulas and
reference code; they are not certificate transcripts.

The actual Python output is stored in \wl{results/python_tests.txt} and
\wl{results/independent_evidence.json}. Native status, novelty, and coverage are
separate files. A missing native transcript is not represented by a fabricated
empty ``success'' report.

\subsection{Proposed source edits}

\wl{patches/apply_focused_patches.py} defaults to a unified-diff dry run. It
requires the pinned Git revision, compares targeted canonical files with
\wl{git show}, requires unique source anchors, and refuses pre-existing backups.
It preflights all selected edits before changing any file. The optional
\wl{--apply} mode writes only the selected canonical sources.

Its scope is deliberately limited: N01 guards the two depth-parser rules;
N02 consumes the canonical backend default; N03 replaces the dependency pattern.
N04 remains a design proposal. The tool does not alter the remote repository,
create a commit, regenerate the root distribution, or claim that a patch has
passed native validation.

The builder edit was actually executed on the local retrieved source fixture.
The patch-anchor safeguards were tested on synthetic anchors for the Wolfram
edits. Those tests do not establish successful loading or correct execution of
a fully patched Wolfram package.

\subsection{Native execution supplied for a later local run}

The nine \wl{focused_regressions.wlt} tests describe desired behavior for N01
and N02. On the unmodified baseline, several are expected to fail. The
fresh-kernel runner loads the local modular package, checks for nine returned
results, and records the report. It does not execute the full upstream suite.
The runner handles a documented modern result-list property and a legacy
alternative, rejecting an unknown report shape rather than counting zero tests
as success.\cite{wltestreport}

The fourteen characterization/comparison cases separately localize N01's
normalizer, core constructor, and public facade; compare held native data for
N04; and exercise selected package/native pairs. Each call is bounded and the
raw expressions and messages are recorded. The reference revision in the output
is explicitly distinguished from a caller-supplied checkout that may contain
proposed edits.

\subsection{Practical sequence}

First reproduce N01 at the private normalizer and core-constructor boundaries,
then at the public facade. Check that the patch changes unsupported complex
exponents into a refusal while retaining real-exponent inputs. Next run the
backend-default controls, including delayed defaults and explicit overrides.
Independently, run the builder mutation tests and its ordinary inlining control.
Finally, inspect N04's actual native return before changing its status policy.

This sequence prioritizes source equality, deterministic API behavior, and the
specific build invariant. It is not another copy of the old general release-CI
roadmap. Existing project validation remains relevant, but no additional native
validation is implied by the Python results in this archive.

\section{Conclusion}

The repository's distinguishing value is the ability to carry a source equation,
a selected real approach, a structured asymptotic representation, and explicit
evidence together. The current implementation has made meaningful progress in
separating those contracts from opaque native delegation. The new audit findings
are correspondingly concentrated at boundaries: a source rewrite that is not an
identity, defaults not consulted during contract selection, a syntactically
incomplete distribution gate, and an over-interpretable native status label.

The immediate mathematical repair is small: require real exponents before
unwrapping the identified logarithms. The deeper constructive result is not to
ban every complex intermediate, but to retain the information it carries when
the overall source is real. The bounded winding coefficient supplies an explicit
example: its correct inverse approximation and uniform cubic error bound can be
written down and independently checked.

The recommendation is therefore to land narrowly evidenced boundary repairs,
validate them with the supplied targeted native probes, and develop extensions
only with explicit source identities and evidence contracts. A longer function
list is less valuable than a reliable explanation of which function a returned
expansion actually approximates.

\clearpage
\appendix
\section{Novelty crosswalk}

\begin{longtable}{@{}p{0.08\textwidth}p{0.19\textwidth}p{0.65\textwidth}@{}}
\toprule
New ID & Closest prior item & Distinct contribution\\
\midrule
\endhead
N01 & C07; C16 & The coefficients become real after an invalid logarithm rewrite; an explicit real monotone source and exact inverse-value sequence disprove the resulting remainder. The issue is source identity, not only codomain admission or opaque Taylor data.\\[5pt]
N02 & D01; C09; B01 & A newly inspected facade declares a configurable backend yet substitutes a literal default instead of consulting its option table. This is distinct from an explicit coefficient option being shadowed by stored model data.\\[5pt]
N03 & V01; V02 & The actual Python assembler accepts eight alternative dependency-call forms. The remedy and tests concern a specific offline build invariant, not generic missing native CI or test inventory.\\[5pt]
N04 & B01; D02 & A concrete all-tree native-call scan cannot distinguish held user data from an unresolved active call. No wrong asymptotic coefficient is asserted.\\
\bottomrule
\end{longtable}

The following inventories were read in addition to the consolidated register.
Their complete pinned paths are in \wl{results/source_coverage.json}.
\begin{longtable}{@{}p{0.12\textwidth}p{0.79\textwidth}@{}}
\toprule
Review & Inventory read\\
\midrule
\endhead
1 & \wl{evidence/findings.csv}\\
2 & README finding table and evidence scope\\
3 & README principal findings and evidence scope\\
4 & \wl{evidence/findings.json}\\
5 & \wl{evidence/findings.json}\\
6 & \wl{evidence/findings.csv}\\
7 & \wl{evidence/findings.csv}\\
8 & \wl{evidence/findings.csv}\\
9 & README finding table and evidence scope\\
10 & \wl{evidence/review_crosswalk.json}\\
11 & \wl{evidence/findings-delta.json}\\
12 & \wl{evidence/novelty_ledger.csv}\\
13 & README detailed principal findings and evidence scope\\
14 & \wl{evidence/findings.json}\\
15 & \wl{evidence/findings.csv}\\
16 & \wl{evidence/findings.csv}\\
17 & \wl{evidence/novelty_matrix.json}\\
18 & \wl{evidence/novelty_ledger.json}\\
\bottomrule
\end{longtable}
This screening does not claim literal full-text coverage of every old article,
annex, and script. Absence from the inventories is supporting evidence of
non-duplication, combined with the mechanism-level comparisons above.

\section{Source-review coverage}

Detailed source inspection covered the ordinary core's public definitions,
exactness checks, model admission, inverse coefficient formulas, symbolic-depth
parser, constructor, object assembly, and selected arithmetic entry points;
the native compatibility module; inverse-function facade plumbing; the interval
certificate implementation; the Gamma forward implementation; the Dirichlet
special-function implementation; the structured native special-function importer;
and substantial special-inverse adapter code.\cite{core,native,inverseexpr,arith,cert,gamma,dirichlet,nativespecial,special}

The root README, kernel module map, compatibility document, review register,
review index, builder source, builder tests, and builder workflow were also read.
Many additional specialized modules were assessed through those documents,
callers, and exported contracts rather than exhaustively line by line. No
uninspected implementation is certified free of defects.

The local executable fixture is only the 4,921-byte upstream Python builder.
The rest of the repository was not reconstructed as an executable checkout.
The supplied source manifest records the distinction, and no full-package
loadability result is inferred from source inspection.

\section{Artifact index}

\begin{longtable}{@{}p{0.46\textwidth}p{0.46\textwidth}@{}}
\toprule
Path & Purpose\\
\midrule
\endhead
\wl{article/asymptotic-current-audit.tex} & Self-contained editable article.\\
\wl{article/asymptotic-current-audit.pdf} & Compiled and visually inspected PDF.\\
\wl{code/independent_checks.py} & Executed mathematics and actual-builder fixture tests.\\
\wl{code/focused_regressions.wlt} & Nine unexecuted desired-contract Wolfram tests.\\
\wl{code/run_native.wls} & Local modular-package test runner.\\
\wl{code/characterize_and_compare.wls} & Fourteen unexecuted native characterization/comparison cases.\\
\wl{code/WindingAwareReference.wl} & Formula-level constructive reference and conditional bound; native-unexecuted.\\
\wl{fixtures/build_standalone_upstream.py} & Exact retrieved upstream builder.\\
\wl{patches/apply_focused_patches.py} & Dry-run-first, pinned-checkout edit proposals for N01--N03.\\
\wl{results/independent_evidence.json} & Actual independent results and mutation outcomes.\\
\wl{results/python_tests.txt} & Actual eleven-test pass log.\\
\wl{results/native_status.json} & Explicit native execution boundary.\\
\wl{results/novelty_ledger.json} & New/old issue crosswalk.\\
\wl{results/source_coverage.json} & Pinned source and prior-inventory coverage.\\
\bottomrule
\end{longtable}

\newpage
\begin{thebibliography}{99}
\small
\bibitem{readme} Vladimir Reshetnikov, \emph{Repository README}. Repository source at the pinned revision.\par
\url{https://github.com/VladimirReshetnikov/Asymptotic/blob/6687962f3c858a4f93623cfc496f33e35c6763d4/README.md}

\bibitem{core} Vladimir Reshetnikov, \emph{Core package: public entry points, models, parseFinite (lines 916--928), construct, and object representation}. Repository source at the pinned revision.\par
\url{https://github.com/VladimirReshetnikov/Asymptotic/blob/6687962f3c858a4f93623cfc496f33e35c6763d4/src/Kernel/AsymptoticAnalysis.wl}

\bibitem{native} Vladimir Reshetnikov, \emph{NativeCompatibility: option table and dispatch (lines 5--80); native-result assembly}. Repository source at the pinned revision.\par
\url{https://github.com/VladimirReshetnikov/Asymptotic/blob/6687962f3c858a4f93623cfc496f33e35c6763d4/src/Kernel/NativeCompatibility.wl}

\bibitem{nativedoc} Vladimir Reshetnikov, \emph{Native compatibility design and limitations}. Repository source at the pinned revision.\par
\url{https://github.com/VladimirReshetnikov/Asymptotic/blob/6687962f3c858a4f93623cfc496f33e35c6763d4/docs/development/NATIVE_COMPATIBILITY.md}

\bibitem{status} Vladimir Reshetnikov, \emph{Consolidated code-review status register}. Repository source at the pinned revision.\par
\url{https://github.com/VladimirReshetnikov/Asymptotic/blob/6687962f3c858a4f93623cfc496f33e35c6763d4/docs/development/CODE_REVIEW_STATUS.md}

\bibitem{index} Vladimir Reshetnikov, \emph{External review index}. Repository source at the pinned revision.\par
\url{https://github.com/VladimirReshetnikov/Asymptotic/blob/6687962f3c858a4f93623cfc496f33e35c6763d4/external-reports/code-review/README.md}

\bibitem{modulemap} Vladimir Reshetnikov, \emph{Kernel module map}. Repository source at the pinned revision.\par
\url{https://github.com/VladimirReshetnikov/Asymptotic/blob/6687962f3c858a4f93623cfc496f33e35c6763d4/src/Kernel/README.md}

\bibitem{inverseexpr} Vladimir Reshetnikov, \emph{InverseFunctionExpressions: public facade and inverse-expression composition}. Repository source at the pinned revision.\par
\url{https://github.com/VladimirReshetnikov/Asymptotic/blob/6687962f3c858a4f93623cfc496f33e35c6763d4/src/Kernel/InverseFunctionExpressions.wl}

\bibitem{arith} Vladimir Reshetnikov, \emph{SeriesArithmetic: arithmetic/evaluation entry points}. Repository source at the pinned revision.\par
\url{https://github.com/VladimirReshetnikov/Asymptotic/blob/6687962f3c858a4f93623cfc496f33e35c6763d4/src/Kernel/SeriesArithmetic.wl}

\bibitem{gamma} Vladimir Reshetnikov, \emph{GammaForward: logarithmic normalization and expansion}. Repository source at the pinned revision.\par
\url{https://github.com/VladimirReshetnikov/Asymptotic/blob/6687962f3c858a4f93623cfc496f33e35c6763d4/src/Kernel/GammaForward.wl}

\bibitem{special} Vladimir Reshetnikov, \emph{SpecialFunctionAdapters: special inverse models and bounds}. Repository source at the pinned revision.\par
\url{https://github.com/VladimirReshetnikov/Asymptotic/blob/6687962f3c858a4f93623cfc496f33e35c6763d4/src/Kernel/SpecialFunctionAdapters.wl}

\bibitem{dirichlet} Vladimir Reshetnikov, \emph{DirichletSpecialFunctions: Zeta and Lerch constructions}. Repository source at the pinned revision.\par
\url{https://github.com/VladimirReshetnikov/Asymptotic/blob/6687962f3c858a4f93623cfc496f33e35c6763d4/src/Kernel/DirichletSpecialFunctions.wl}

\bibitem{nativespecial} Vladimir Reshetnikov, \emph{NativeSpecialFunctions: structured native import}. Repository source at the pinned revision.\par
\url{https://github.com/VladimirReshetnikov/Asymptotic/blob/6687962f3c858a4f93623cfc496f33e35c6763d4/src/Kernel/NativeSpecialFunctions.wl}

\bibitem{cert} Vladimir Reshetnikov, \emph{InverseCertificates: interval primitives and adaptive certification}. Repository source at the pinned revision.\par
\url{https://github.com/VladimirReshetnikov/Asymptotic/blob/6687962f3c858a4f93623cfc496f33e35c6763d4/src/Kernel/InverseCertificates.wl}

\bibitem{builder} Vladimir Reshetnikov, \emph{Standalone package builder}. Repository source at the pinned revision.\par
\url{https://github.com/VladimirReshetnikov/Asymptotic/blob/6687962f3c858a4f93623cfc496f33e35c6763d4/validation/build_standalone.py}

\bibitem{buildtests} Vladimir Reshetnikov, \emph{Standalone builder tests}. Repository source at the pinned revision.\par
\url{https://github.com/VladimirReshetnikov/Asymptotic/blob/6687962f3c858a4f93623cfc496f33e35c6763d4/validation/test_standalone.py}

\bibitem{workflow} Vladimir Reshetnikov, \emph{Standalone freshness workflow}. Repository source at the pinned revision.\par
\url{https://github.com/VladimirReshetnikov/Asymptotic/blob/6687962f3c858a4f93623cfc496f33e35c6763d4/.github/workflows/standalone.yml}

\bibitem{r11} Vladimir Reshetnikov, \emph{Review 11: finding-delta ledger}. Repository source at the pinned revision.\par
\url{https://github.com/VladimirReshetnikov/Asymptotic/blob/6687962f3c858a4f93623cfc496f33e35c6763d4/external-reports/code-review/wave-2/code-review-11/evidence/findings-delta.json}

\bibitem{r17} Vladimir Reshetnikov, \emph{Review 17: novelty matrix, including prior native Zeta overlap}. Repository source at the pinned revision.\par
\url{https://github.com/VladimirReshetnikov/Asymptotic/blob/6687962f3c858a4f93623cfc496f33e35c6763d4/external-reports/code-review/wave-2/code-review-17/evidence/novelty_matrix.json}

\bibitem{r18} Vladimir Reshetnikov, \emph{Review 18: novelty ledger}. Repository source at the pinned revision.\par
\url{https://github.com/VladimirReshetnikov/Asymptotic/blob/6687962f3c858a4f93623cfc496f33e35c6763d4/external-reports/code-review/wave-2/code-review-18/evidence/novelty_ledger.json}

\bibitem{wlseries} Wolfram Research, \emph{Series}, Wolfram Language Documentation. Accessed 9 September 2026.\par
\url{https://reference.wolfram.com/language/ref/Series.html}

\bibitem{wlseriesdata} Wolfram Research, \emph{SeriesData}, Wolfram Language Documentation. Accessed 9 September 2026.\par
\url{https://reference.wolfram.com/language/ref/SeriesData.html}

\bibitem{wlinverseseries} Wolfram Research, \emph{InverseSeries}, Wolfram Language Documentation. Accessed 9 September 2026.\par
\url{https://reference.wolfram.com/language/ref/InverseSeries.html}

\bibitem{wlcompose} Wolfram Research, \emph{ComposeSeries}, Wolfram Language Documentation. Accessed 9 September 2026.\par
\url{https://reference.wolfram.com/language/ref/ComposeSeries.html}

\bibitem{wlasymptotic} Wolfram Research, \emph{Asymptotic}, Wolfram Language Documentation. Accessed 9 September 2026.\par
\url{https://reference.wolfram.com/language/ref/Asymptotic.html}

\bibitem{wlasymsolve} Wolfram Research, \emph{AsymptoticSolve}, Wolfram Language Documentation. Accessed 9 September 2026.\par
\url{https://reference.wolfram.com/language/ref/AsymptoticSolve.html}

\bibitem{wlasymintegrate} Wolfram Research, \emph{AsymptoticIntegrate}, Wolfram Language Documentation. Accessed 9 September 2026.\par
\url{https://reference.wolfram.com/language/ref/AsymptoticIntegrate.html}

\bibitem{wlasymdsolve} Wolfram Research, \emph{AsymptoticDSolveValue}, Wolfram Language Documentation. Accessed 9 September 2026.\par
\url{https://reference.wolfram.com/language/ref/AsymptoticDSolveValue.html}

\bibitem{wlpowerexpand} Wolfram Research, \emph{PowerExpand}, Wolfram Language Documentation. Accessed 9 September 2026.\par
\url{https://reference.wolfram.com/language/ref/PowerExpand.html}

\bibitem{wllog} Wolfram Research, \emph{Log}, Wolfram Language Documentation. Accessed 9 September 2026.\par
\url{https://reference.wolfram.com/language/ref/Log.html}

\bibitem{wlsetoptions} Wolfram Research, \emph{SetOptions}, Wolfram Language Documentation. Accessed 9 September 2026.\par
\url{https://reference.wolfram.com/language/ref/SetOptions.html}

\bibitem{wlexpressions} Wolfram Research, \emph{Expressions: special ways to input expressions}, Wolfram Language Documentation. Accessed 9 September 2026.\par
\url{https://reference.wolfram.com/language/tutorial/Expressions.html}

\bibitem{wlapply} Wolfram Research, \emph{Apply}, Wolfram Language Documentation. Accessed 9 September 2026.\par
\url{https://reference.wolfram.com/language/ref/Apply.html}

\bibitem{wltestreport} Wolfram Research, \emph{TestReport}, Wolfram Language Documentation. Accessed 9 September 2026.\par
\url{https://reference.wolfram.com/language/ref/TestReport.html}

\end{thebibliography}
\end{document}
